\begin{pastexam}[label={px:ntu-108-ee-dsb-15}]{臺大電機所 108 學年度（B）第 15 題}
（複選；每題 10 分，每個選項單獨計分，答對 2 分，答錯倒扣 2 分）

Which of the following statements are true?
\begin{enumerate}[label=(\Alph*).]
\item Hash tables are generally used in cases where insertion and deletion operations are more
often than searches.
\item Open-addressed hash tables cannot be used if the data does not have unique keys.
\item With linear probing, searches become faster as we remove data from the hash table.
\item Both linear and quadratic probing can probe all locations in the table.
\item Double hashing resolves the clustering issue by using key-dependent probe intervals.
\end{enumerate}
\end{pastexam}

\begin{pastsoln}
答案是 (E)。

\textbf{(A) 錯}。 雜湊表的長處是用鍵直接算出位置，\textbf{搜尋}期望 $O(1)$，所以它用在「查詢
很頻繁」的場合，例如符號表、字典、快取。 如果只是大量插入、刪除而很少搜尋，用不著雜湊：一個
未排序的串列插入就是 $O(1)$。 而且雜湊表的插入、刪除本身也要先搜尋（插入要確認鍵不重複、刪除
要先找到它），搜尋快是一切的基礎。

\textbf{(B) 錯}。 開放定址的探測機制沒有要求鍵不重複。 允許重複鍵時，插入就沿探測序列放進第一個
空格（不必先檢查有沒有相同的鍵）；要找出某個鍵的所有紀錄，就沿探測序列一直走到空格為止，把
沿途鍵相同的都收集起來。 這和鏈結法要掃完整條鏈才能找齊重複鍵是同一回事。 常見的實作（例如
資結書虛擬碼~12.3 的 \proc{Mark-Insert}）遇到相同的鍵就不再插入，那是「集合」這種用法的選擇，不是
開放定址做不到。

\textbf{(C) 錯}。 開放定址的刪除不能把格子清空，否則會切斷經過那一格的探測序列，要改成放刪除
標記（資結書第~12.3.4~節）。 搜尋遇到刪除標記時必須當成有東西、繼續往下找，所以刪除之後探測
序列並沒有變短：失敗的搜尋照樣要走到真正的空格才停。 刪除標記越積越多，搜尋反而越來越慢，
這也是用刪除標記的表要定期重建的原因。 （線性探測另有「刪除後把後面的鍵往前搬」的做法，可以真的
縮短序列，但那不是開放定址的一般刪除方式，題目的敘述不能當成一般性質。）

\textbf{(D) 錯}。 線性探測 $h(k) + i$ 會走遍所有 $m$ 格。 二次探測 $h(k) + i^2$ 不一定：$m$ 是
奇質數時只保證前 $(m + 1)/2$ 個位置互不相同（資結書定理~12.3），而 $i^2 \bmod m$ 只有
$(m+1)/2$ 種值，所以最多只探得到一半多一點的格子。 例如 $m = 7$：$i^2 \bmod 7$ 只有
$0, 1, 4, 2$ 四種，表有 $7$ 格卻只探得到 $4$ 格，其他 $3$ 格都空著時插入也可能失敗。

\textbf{(E) 對}。 雙重雜湊 $h(k, i) = (h_1(k) + i\, h_2(k)) \bmod m$ 的間隔 $h_2(k)$ 由鍵決定
（資結書第~12.3~節）。 兩個鍵即使 $h_1$ 相同，間隔通常不同，探測序列很快分開；線性探測的主要
叢集（連成一大片）和二次探測的次要叢集（同 $h(k)$ 的鍵序列完全相同）都因此大為減輕。
\end{pastsoln}
